%% kapa: exam=1 round=26 year=2015 period=1 form=A subject=civil_law subject_ko=민법 seq=36 number=36
%% source: Q-Net 감정평가사 기출문제 (cst00309) artlSeq=5138001 file=extracted/1008/5138001/제26회_감정평가사 1차_기출/제26회_감정평가사 1차_기출/1교시_감평_A형.pdf
%% conversion: kordoc 4.24.51 -> tools/gen1.py (auto); review=proofread-sonnet; flags=none
\begin{question}{36}
법정지상권에 관한 설명으로 옳지 않은 것은? (다툼이 있으면 판례에 따름)
\begin{choices}
  \item 대지와 건물이 동일한 소유자에 속한 경우에 건물에 전세권을 설정한 때에는 그 대지소유권의 특별 승계인은 전세권설정자에 대하여 지상권을 설정한 것으로 본다.
  \item 토지공유자 중 1인이 다른 공유자 지분 과반수의 동의를 얻어 건물을 건축한 후 토지와 건물의 소유자가 달라진 경우, 관습상의 법정지상권이 성립된다.
  \item 미등기건물을 그 대지와 함께 매도하여 대금이 완납되었으나 건물이 미등기인 관계로 대지에 관하여만 매수인 앞으로 소유권이전등기가 경료된 경우, 매도인에게 관습상의 법정지상권은 인정되지 않는다.
  \item 관습상 법정지상권이 성립하려면 토지와 그 지상 건물이 애초부터 동일인의 소유에 속하였을 필요는 없고, 그 소유권이 유효하게 변동될 당시에 동일인이 토지와 그 지상 건물을 소유하였던 것으로 족하다.
  \item 토지와 건물의 소유자 甲으로부터 乙이 건물을 매수하여 취득한 경우 乙이 건물 소유를 위해 甲과 대지의 임대차계약을 체결하였다면, 관습상 법정지상권을 포기한 것으로 본다.
\end{choices}
\end{question}
\answer{2}
